\documentclass[11pt,a4paper]{article}

\usepackage[utf8]{inputenc}
\usepackage[T1]{fontenc}
\usepackage[ngerman]{babel}
\usepackage{amsmath,amssymb,amsthm}
\usepackage{booktabs}
\usepackage{hyperref}

\title{Ein komplexes EABC-Primvierlingsfeld und seine Vierer-Phasenrotation}
\author{Thomas Hoffbauer}
\date{\today}

\newtheorem{definition}{Definition}
\newtheorem{hypothesis}{Hypothese}
\newtheorem{observation}{Beobachtung}

\begin{document}
	
	\maketitle
	
	\begin{abstract}
		Wir untersuchen Primvierlinge der Form
		\[
		(p,p+2,p+6,p+8)
		\]
		im Rahmen einer EABC-Klassifikation modulo \(12\). Dabei werden die beiden beobachteten Umlaufphasen \(ABCE\) und \(CEAB\) zu einem komplexen Primvierlingsfeld kombiniert. Numerische Tests bis \(10^7\) zeigen, dass die Amplitude dieses Feldes skalenabhängig mit einem aus Riemann-Nullstellen konstruierten Zeta-Detektor korreliert, während die Phase selbst zeta-schwach bleibt. Die Phasendifferenz besitzt jedoch ein starkes internes Spektrum mit dominanter Periode \(4\). Dies legt nahe, dass die Zeta-Funktion primär die Dichte-Amplitude sieht, während die EABC-Struktur eine eigenständige Phasenrotation trägt.
	\end{abstract}
	
	\section{EABC-Klassifikation}
	
	Wir betrachten die vier primen Restklassen modulo \(12\):
	\[
	E\equiv 1,\qquad
	A\equiv 5,\qquad
	B\equiv 7,\qquad
	C\equiv 11
	\pmod{12}.
	\]
	
	Diese Klassen bilden mit der Multiplikation modulo \(12\) eine Kleinsche Viererstruktur:
	\[
	A^2=B^2=C^2=E,
	\qquad
	AB=C,\quad AC=B,\quad BC=A.
	\]
	
	\section{Primvierlinge und Umlaufphasen}
	
	Ein Primvierling besitzt die Form
	\[
	(p,p+2,p+6,p+8).
	\]
	
	Für geeignete Mittelpunkte
	\[
	m=p+4
	\]
	erscheinen zwei Phasenlagen:
	\[
	ABCE
	\quad\text{und}\quad
	CEAB.
	\]
	
	Diese beiden Folgen sind keine entgegengesetzten Chiralitäten, sondern zyklische Phasenlagen desselben orientierten Umlaufs.
	
	\begin{definition}[Band- und Phasenkanal]
		Seien \(Y_{ABCE}(m)\) und \(Y_{CEAB}(m)\) die Indikatorfelder der beiden Vierlingsphasen. Wir definieren
		\[
		Y_{\mathrm{band}}(m)
		=
		Y_{ABCE}(m)+Y_{CEAB}(m)
		\]
		und
		\[
		Y_{\mathrm{phase}}(m)
		=
		Y_{ABCE}(m)-Y_{CEAB}(m).
		\]
	\end{definition}
	
	\section{Komplexes Primvierlingsfeld}
	
	\begin{definition}[Komplexes EABC-Feld]
		Das komplexe Primvierlingsfeld sei
		\[
		\Psi(m)
		=
		Y_{\mathrm{band}}(m)
		+
		iY_{\mathrm{phase}}(m).
		\]
		Dann ist
		\[
		|\Psi(m)|
		\]
		der Dichte- beziehungsweise Amplitudenkanal, während
		\[
		\arg \Psi(m)
		\]
		den EABC-Phasenkanal beschreibt.
	\end{definition}
	
	\section{Holonomieoperator}
	
	Wir definieren den Vorwärtsumlauf
	\[
	P_+:\quad A\to B\to C\to E\to A
	\]
	und den Rückwärtsumlauf
	\[
	P_-:\quad A\to E\to C\to B\to A.
	\]
	
	Daraus entstehen der Bandoperator
	\[
	H_{\mathrm{band}}=P_+ + P_-
	\]
	und der chirale beziehungsweise gerichtete Operator
	\[
	H_{\mathrm{chiral}}=P_+ - P_-.
	\]
	
	Numerisch ergibt sich
	\[
	\operatorname{spec}(H_{\mathrm{chiral}})
	=
	\{0,0,2i,-2i\}.
	\]
	
	Damit ist \(H_{\mathrm{chiral}}\) ein echter Rotationsoperator des EABC-Phasenraums.
	
	\section{Numerische Beobachtungen}
	
	Die Auswertung bis \(10^7\) ergab:
	
	\begin{center}
		\begin{tabular}{l r}
			\toprule
			Struktur & Anzahl bis \(10^7\) \\
			\midrule
			Zwilling links & 19797 \\
			Zwilling Mitte & 19542 \\
			Zwilling rechts & 19674 \\
			Drilling links & 4319 \\
			Drilling rechts & 4305 \\
			Primvierling & 898 \\
			Fünfling I & 160 \\
			Fünfling II & 161 \\
			Sechsling & 18 \\
			\bottomrule
		\end{tabular}
	\end{center}
	
	Der Bandkanal zeigt gegen den Zeta-Detektor eine robuste skalenabhängige Korrelation. Für kleine Bin-Größen wurden signifikante Werte gefunden, etwa:
	\[
	\rho_{B=100}\approx -0.241,
	\qquad
	\rho_{B=1000}\approx -0.193.
	\]
	
	Shuffle-Tests bestätigen diesen Befund für kleine bis mittlere Bins.
	
	Der Phasenkanal dagegen bleibt gegen Zeta weitgehend schwach:
	\[
	\operatorname{corr}(Y_{\mathrm{phase}},Z)\approx 0.00047.
	\]
	
	\section{Viererperiode der Phase}
	
	Die Phasenänderung
	\[
	d\phi(m)=\phi(m+1)-\phi(m)
	\]
	mit
	\[
	\phi(m)=\arg\Psi(m)
	\]
	zeigt ein starkes internes Spektrum.
	
	Die Autokorrelation ergab unter anderem:
	\[
	\mathrm{ACF}(2)\approx -0.691,
	\qquad
	\mathrm{ACF}(4)\approx +0.691,
	\qquad
	\mathrm{ACF}(6)\approx -0.694,
	\qquad
	\mathrm{ACF}(8)\approx +0.691.
	\]
	
	Das Leistungsspektrum besitzt eine dominante Frequenz
	\[
	f=0.25,
	\]
	also eine Periode
	\[
	T=4.
	\]
	
	\begin{observation}[EABC-Phasenrotation]
		Der Phasenkanal des komplexen Primvierlingsfeldes besitzt eine dominante Viererperiode. Diese Periode entspricht exakt der EABC-Struktur.
	\end{observation}
	
	\section{Interpretation}
	
	Die numerischen Ergebnisse legen folgende Trennung nahe:
	
	\[
	\boxed{
		\zeta \text{ koppelt an } |\Psi|
	}
	\]
	
	und
	
	\[
	\boxed{
		H_{\mathrm{chiral}} \text{ erzeugt } \arg \Psi.
	}
	\]
	
	Die Riemann-Zeta-Funktion wirkt damit als Detektor der skalaren Primdichte, während die EABC-Struktur einen eigenständigen Phasenraum beschreibt.
	
	\begin{hypothesis}[Komplexe Primvierlingsfeld-Hypothese]
		Die Primvierlingsstruktur besitzt ein komplexes Feld
		\[
		\Psi=Y_{\mathrm{band}}+iY_{\mathrm{phase}},
		\]
		dessen Amplitudenanteil zeta-sensitiv ist, während sein Phasenanteil eine interne EABC-Holonomie mit dominanter Viererrotation trägt.
	\end{hypothesis}
	
	\section{Ausblick}
	
	Der nächste Schritt besteht darin, den Phasenkanal nicht gegen die Zeta-Nullstellen, sondern gegen das Eigenwertspektrum eines eigenständigen EABC-Holonomieoperators zu testen. Insbesondere sollten Autokorrelationen, Fourier-Spektren und Übergangsmatrizen des Phasenfeldes bis \(10^8\) und \(10^9\) untersucht werden.
	
\end{document}